\documentclass[11pt]{article}
\usepackage[margin=1in]{geometry}
\usepackage{amsmath, amssymb}
\usepackage{hyperref}
\usepackage{enumitem}
\setlist[itemize]{noitemsep,topsep=2pt}

\title{Memo: Reframing the Outer Sorting Model}
\author{}
\date{March 07, 2026}

\begin{document}
\maketitle

\section{Main point}
The current outer model asks filing rates to do too much. Filings are simultaneously:
\begin{itemize}
    \item the main sorting signal tenants condition on,
    \item a key outcome the model must match,
    \item and a channel through which landlord type and tenant default risk enter equilibrium.
\end{itemize}
This is convenient computationally, but it is likely too strong as a model of tenant information. In practice, tenants are much more likely to observe or infer \emph{building characteristics} than filing rates. Observable cues such as building age, visible condition, size, type, block quality, and neighborhood characteristics are likely to drive most sorting, while filing behavior remains only imperfectly observed.

\section{Implication for the signal structure}
The outer model should lean less on filings as the sorting signal. A cleaner information structure is:
\[
x_b^{\text{sort}} = x_b^{\text{obs}} + x_b^{\text{latent}},
\]
where:
\begin{itemize}
    \item $x_b^{\text{obs}}$ depends on observable building characteristics,
    \item $x_b^{\text{latent}}$ captures landlord aggressiveness / harshness / under-maintenance risk that tenants do not observe directly.
\end{itemize}

The key distinction is:
\begin{itemize}
    \item tenants sort primarily on $x_b^{\text{obs}}$ (possibly with noise),
    \item the econometrician uses filings, complaints, and maintenance outcomes to infer the latent harshness component ex post.
\end{itemize}

This makes the model's information assumptions more realistic and avoids giving tenants too much access to the same ``badness'' object used by the econometrician.

\section{Revised outer sorting block}
Instead of
\[
s_b^\star(\xi)=\Lambda\!\left(\xi + \Delta\alpha\,\bar f_b\right),
\]
the sorting rule should be re-centered around observable building characteristics:
\[
s_b^\star(\xi)=\Lambda\!\left(\xi + \Delta_X X_b + \nu_b\right),
\]
or more generally
\[
s_b^\star(\xi)=\Lambda\!\left(\xi + \Delta_X X_b + \Delta_q q_b\right),
\]
where:
\begin{itemize}
    \item $X_b$ is a vector of observable building characteristics,
    \item $\nu_b$ is a reduced-form latent attractiveness or harshness shifter,
    \item $q_b$ is a latent building-quality / landlord-harshness index if one wants a more structural representation.
\end{itemize}

The citywide mass constraint remains:
\[
\sum_b U_b s_b = \mu \sum_b U_b.
\]

\paragraph{Interpretation.}
This specification says that segregation emerges because constrained tenants sort differently across observable building environments. It does \emph{not} require that tenants observe landlord type or filing propensity directly.

\section{Where filings belong}
Filings should remain an \emph{equilibrium outcome}, not the main sorting signal. The natural role for filings in the model is:
\begin{itemize}
    \item landlord type $\theta_b$ raises filing propensity,
    \item expected tenant default risk raises filing propensity,
    \item complaints may later feed into filings, if that channel is added.
\end{itemize}

A natural filing block is therefore:
\[
\bar \delta_b = s_b\delta_H + (1-s_b)\delta_L,
\]
\[
\log \bar f_b = a_F + b_{F\theta}\mathbf 1\{\theta_b=B\} + b_{F\delta}\bar\delta_b,
\]
or, if needed in the data-generating process,
\[
\log \bar f_b = a_F + b_{F\theta}\mathbf 1\{\theta_b=B\} + b_{F\delta}\bar\delta_b + b_{FC}\log(1+\bar c_b).
\]

In this formulation:
\begin{itemize}
    \item tenant-side risk enters filings through $\bar\delta_b$,
    \item landlord type enters filings through $\theta_b$,
    \item filings are generated by the environment rather than driving the entire sorting structure.
\end{itemize}

\section{Why this is preferable}
This revision has three advantages.

\subsection*{1. More realistic tenant information}
Tenants are more likely to observe building characteristics than filing rates. The revised model reflects this.

\subsection*{2. Cleaner separation between sorting and outcomes}
Sorting is driven by observable or weakly latent building environment, while filings are one outcome of landlord behavior and tenant risk. This prevents filings from doing too much work.

\subsection*{3. Better discipline for calibration}
If the model uses filings as both the signal and the outcome, then matching filing dispersion can distort the sorting moments. Separating these objects should make calibration more stable and interpretation cleaner.

\section{Tradeoff}
The cost of this revision is that identification becomes harder. Once sorting is not driven mainly by filings, the model needs sufficient variation in:
\begin{itemize}
    \item observable building characteristics,
    \item latent harshness/quality shifters,
    \item landlord type,
    \item and tenant default risk
\end{itemize}
to match sorting, complaint, and filing moments jointly.

This is a worthwhile tradeoff. The model becomes more realistic and the interpretation of filings becomes cleaner.

\section{Role of three landlord types}
Moving from two landlord types to three may eventually help fit richer distributional shapes, especially for filing rates and, later, price curves. But that is not the immediate next fix. The higher priority change is to revise the \emph{sorting signal}. Once the sorting and filing blocks are stable, a three-type extension can be used to improve distributional fit in a disciplined way.

\section{Recommended next step}
The next revision of the outer model should:
\begin{itemize}
    \item make the sorting signal depend primarily on observable building characteristics,
    \item keep filings as an equilibrium outcome driven by landlord type and average default risk,
    \item and leave three landlord types for a later extension if the two-type model still underfits the observed distributions.
\end{itemize}

A concise way to summarize the intended model is:
\begin{quote}
Tenants sort on noisy signals of observable building quality, while filings are an equilibrium outcome driven by landlord type and tenant default risk.
\end{quote}

\end{document}
